% theme: science_and_technology
\MajorQuestion{科学技術}
\SetRowSpaceQanda

\Instruction{次の問いに答えなさい。}
\QQ{石油などを原料として人工的につくられる有機物の素材を\NamedBlank{plastic}という。}{プラスチック}
\QQ{$\alpha$線、$\beta$線、$\gamma$線、X線、中性子線などをまとめて\NamedBlank{radiation}という。}{放射線}
\QQ{\RefBlank{radiation}を出す能力を\NamedBlank{radioactivity}という。}{放射能}
\QQ{太陽光や風力など、使い続けてもなくならないエネルギーを\NamedBlank{renewable_energy}という。}{再生可能エネルギー}
\QQ{コンピュータが学習や判断などを行う技術を、アルファベット2文字で何というか。}{AI}

\Instruction{\strut}
\QQ{\RefBlank{plastic}に含まれる主な元素は何か、2つ答えなさい。}{炭素、水素}
\QQ{\RefBlank{radioactivity}の大きさを表す単位を何というか。}{ベクレル（Bq）}
\QQ{高い所から流れ落ちる水で水車を回す発電方法を何というか。}{水力発電}
\QQ{資源の消費を減らし、再利用を進めて、資源を循環させる社会を何というか。}{循環型社会}
\QQ{自然界へ流れ出た\RefBlank{plastic}が細かく砕かれ、粒子になったものを何というか。}{マイクロプラスチック}

\Instruction{\strut}
\QQ{\RefBlank{radioactivity}をもつ物質を何というか。}{放射性物質}
\QQ{化石燃料を燃やして高温・高圧の水蒸気をつくり、タービンを回す発電方法を何というか。}{火力発電}
\QQ{発電時に生じる排熱の一部を、給湯や暖房にも利用するしくみを何というか。}{コージェネレーションシステム}
\QQ{\RefBlank{radiation}が人体に与える影響の大きさを表す単位を何というか。}{シーベルト（Sv）}
\QQ{特定のすぐれた性質やはたらきをもつようにつくられた高分子を何というか。}{機能性高分子}

\Instruction{\strut}
\QQ{硫黄酸化物や窒素酸化物などの影響で、通常より酸性が強くなった雨を何というか。}{酸性雨}
\QQ{土壌中の微生物によって分解される\RefBlank{plastic}を何というか。}{生分解性プラスチック}
\QQ{原子核が分裂するときに出るエネルギーを利用して、水蒸気でタービンを回す発電方法を何というか。}{原子力発電}
\QQ{生活に必要なものやエネルギーを、現在だけでなく将来の世代も安定して得られる社会を何というか。}{持続可能な社会}
\QQ{大気中の二酸化炭素濃度の上昇などにより、地球の平均気温が上がる現象を何というか。}{地球温暖化}
